%Superposition
\newpage

%Einleitung
\newvideofile{Superpos}{Superpositionsprinzip}
\begin{frame}
	\fta{Principle of superposition}

	\s{ 
        After a basic analysis of meshes and nodes in electrical 
        networks, increasingly complex electrical networks are considered.
        For this purpose, the superposition method according to Helmholtz is used. 
        The superposition method is also referred to as the principle of superposition.     
        In the superposition principle, all sources are evaluated individually one after the other 
        and the results of the individual calculations are superimposed.
    }


    \begin{Learning objectives}{Principle of superposition}
        The students 
        \begin{itemize}
            \item understand the conditions of systems theory for the analysis of direct current networks. 
            \item can apply the superposition principle to electrical networks. 
        \end{itemize}                   
    \end{Learning objectives}

    \speech{SuperposLernz}{1}{Das Superpositionsprinzip: ,,,,
    Verfügt ein elektrisches Netzwerk über mehr als eine Quelle, werden die Komponenten 
    durch die Quellen gegenseitig überlagert. Diese Überlagerung wird durch das 
    Superpositionsprinzip für jede Quelle aufgelöst analysiert. Hierfür werden zuvor 
    ein paar Grundlagen für das Superpositionsprinzip aus der Systemtheorie erläutert, 
    um anschließend den Überlagerungssatz zu erklären.} 
\end{frame}

%Exkurs Systemtheorie
\begin{frame}
    \renewcommand{\figurename}{Figure}
    \ftb{Digression Systems theory}

    \s{
        If only the input signal and the output signal are considered in a model, this is referred to as
        a black box. Here, no information is given about what happens between these 
        two signals. Figure \ref{BildBlackBox} describes such a model. Here, 
        a black box is considered, which has an input variable $F_\mathrm{E}$ and an output variable $F_\mathrm{A}$.
        In the black box, the input signal is transformed into the output signal. 
        The image as a whole is referred to as a system. 
    }

    \onslide<1->{
    \b{
    \begin{itemize}    
		\item Betrachtung einer Blackbox als System mit Eingangsgrößen und Ausgangsgrößen.
    \end{itemize}
    }

    \fu{
        \input{Tikz/src/blackbox.tex}
        }{{\bf Black box in systems theory.} A black box with input variables on the left and output variables on the right.
        \label{BildBlackBox}}}

    \onslide<2->{
    \b{
    \begin{itemize}    
		\item Transformation des Eingangssignals $F_\mathrm{E}$ in ein Ausgangssignal $F_\mathrm{A}$.
    \end{itemize}
    }

    \s{
        Based on the black box presented, the output signal $F_\mathrm{A}$ can be determined as a function 
        of the input signal $F_\mathrm{E}$. This transformation can be described by equation 
        \ref{GleichungFunktionSystem}.
    }

    \begin{eq}
        F_\mathrm{A}=T(F_\mathrm{E})      \label{GleichungFunktionSystem}
    \end{eq}}

    \speech{SuperposExkurs}{1}{Es wird ein System mit einer Eingangsgröße F E und einer Ausgangsgröße F A dargestellt.
    Die Ausgangsgröße ergibt sich dabei aus der Transformation T der Eingangsgröße. In der Systemtheorie werden Systeme, 
    bei denen keine Informationen über die Transformation selber angegeben werden auch Black Box gegannt.}
    \speech{SuperposExkurs}{2}{Hierbei kann die Transformation T der Eingangsgröße zur Ausgangsgröße beispielsweise durch die 
    Funktion einer Spannungsquelle oder Stromquelle erfolgen.}
\end{frame}

\begin{frame}
    \renewcommand{\figurename}{Figure}
    \ftx{Black Box}

    \s{
        An electrical network can also be considered such a system. The electrical network shown in Figure
        \ref{BildElektrischesSystem} has an input voltage 
        $U_\mathrm{E}$ as the input variable and an output voltage $U_\mathrm{A}$ as the output variable. The black box, which 
        represents the system to be transformed, consists of the two resistors $R_1$ and $R_2$.
  
    }

    \onslide<1->{
    \b{
        \begin{itemize}    
            \item Beschreibung einer Black Box als elektrischen Netzwerk mit einer Eingangsspannung
            und einer Ausgangsspannung.
        \end{itemize}
    }}

    \onslide<2->{
    \fu{
        \input{Tikz/src/spannungsteilerblackbox.tex}
        
        }{{\bf Voltage divider as a black box.} Consideration of an electrical network as a black box with input voltage and output voltage.
        \label{BildElektrischesSystem}}}

    \s{
        The relationship between the output signal of a system and the input signal from Equation \ref{GleichungFunktionSystem} still applies. The output signal 
        can be described as a transformation of the input signal as in Equation \ref{GleichungNetzwerk}. 
        The output signal is equal to the voltage drop across resistor $R_2$. 
        This voltage can be calculated from the product of the input voltage and the resistance ratio.
        This relationship was also discussed in the section on voltage dividers in Chapter 4.    
    }

    \onslide<3->{
    \b{
        \begin{itemize}    
            \item Transformation der Eingangsspannung $U_\mathrm{E}$ in eine Ausgangsspannung $U_\mathrm{A}$.
        \end{itemize}
    }

    \begin{eq}
        U_\mathrm{A}=T(U_\mathrm{E})=U_\mathrm{E}\cdot\frac{R_2}{R_1+R_2}  \label{GleichungNetzwerk}
    \end{eq}}

    \speech{SuperposBlackBox}{1}{Wird die Black Box näher betrachtet, offenbart sie das zu transformierende System. Angenommen:
    das zu Transformierende System sei ein elektrisches Netzwerk, so könnte dieses wie folgt dargestellt werden.}
    \speech{SuperposBlackBox}{2}{Das zu Transformierende System ist hier ein einfacher Spannungsteiler aus den beiden Widerständen
    R 1 und R 2. Die Eingangsspannung U E des Spannungsteilers liegt über der Serienschaltung an. Die Ausgangsspannung U A liegt über 
    den Widerstand R 2 an.}
    \speech{SuperposBlackBox}{3}{Die Ausgangsspannung kann so als transformiertes System in Abhängigkeit von der Eingangsspannung 
    definiert werden. So wird in diesem Beispiel die Eingangsspannung mit dem Spannungsteilerverhältnis multipliziert, um die Ausgangsspannung
    ermitteln zu können.}
\end{frame}

\begin{frame}
    \renewcommand{\figurename}{Figure}
    \ftx{LTI-Systeme}

    \s{
        \textbf{Causality}\\

        If an output variable is based exclusively on the transformation of an input variable, there is
        a direct relationship between cause and effect. This also means that the output variable does not provide a changing system response prior to 
        an excitation. A system in which these conditions prevail is referred to as a causal system (see Figure \ref{BildKausalesSystem}).
        For example, a DC excitation of a system produces a time-indifferent
        DC response. A system that behaves differently is referred to as a non-causal system. 
        If the value of the input signal is zero for $t < t_0$, the value of the output signal 
        must also be zero for the same period of time in order to fulfil the conditions of a causal system. 
    }

    \b{
        \begin{itemize}    
            \item<1-> Kausale Systeme zeigen keine Sytemantwort vor der Anregung
            \item<2-> Zeitinvariante Systeme reagieren immer gleich auf ein Eingangssignal, unabhängig vom Zeitpunkt
            \item<2-> Lineares System $\rightarrow$ endliche Ausgangsamplitude bei endlicher Eingangsamplitude
            \item<3-> Zeitinvariante und lineare System werden als LTI-Systeme kategorisiert
        \end{itemize}
    }

    \onslide<2->{
    \fu{
        \resizebox{.9\textwidth}{!}{%
        \begin{minipage}{0.2\textwidth}
            \input{Tikz/src/transformation1.tex}	
		\end{minipage}
        \hspace{0.1cm}
        \begin{minipage}{0.5\textwidth}
            \input{Tikz/src/transformation2.tex}	
		\end{minipage}
        \hspace{0.1cm}
        \begin{minipage}{0.2\textwidth}
            \input{Tikz/src/transformation3.tex}	
		\end{minipage}
        }
	}{{\bf Transformation of a causal system.} System transformed from $x(t)$ to $y(t)$ for a causal system. \label{BildKausalesSystem}}}

    \s{
        \textbf{Time variance}\\
        
        If the output signal of a system always responds to an input signal in the same way over time, 
        this is referred to as a time-invariant system. If the timing of the input signal shifts, for example 
        by $t_0$, the output signal must still respond in the same way, based on the input signal, in time-invariant systems.
        
        The stability of a system also describes a relationship between the input signal
        and the output signal. If an input signal with finite amplitude produces an output signal that does not grow beyond
        all limits, the system is stable. If the output signal continues to grow after 
        the input signal has been deactivated, the system is unstable. 

        \begin{eq}
            T(\alpha \cdot u_\mathrm{e})=\alpha \cdot T(u_\mathrm{e})    \label{GleichungLinearität}
        \end{eq}

        \textbf{Linearity}\\

        If the response of the output signal is proportional to the excitation of the input signal, the
        system is described as a \textbf{linear} system (Equation \ref{GleichungLinearität}). If two superimposed 
        input signals act on a system, they can be considered separately and added together in this case. In this way, 
        the transformations of the input signals, as in Equation \ref{GleichungÜberlagerteSignale},
        can also be considered separately.
 

        \begin{eq}
            T(u_1+u_2)=T(u_1)+T(u_2)    \label{GleichungÜberlagerteSignale}
        \end{eq}

        Systems that are both time-invariant and linear are referred to as LTI (Linear Time Invariant) systems.
        The systems in this module are approximated as LTI systems. Components
        that exhibit non-linearities are dealt with, for example, in the chapter on periodic quantities. 
    }   

    \speech{SuperposLTI}{1}{Ein System lässt sich durch verschiedene Kriterien beschreiben. So zeigen kausale Systeme vor der Anregung keine Systemantwort.}
    \speech{SuperposLTI}{2}{Zeitinvariante Systeme reagieren immer gleich auf ein Eingangssignal, dabei ist es unabhängig ob der Zeitpunkt T null gewählt wird
    oder ein früher beziehungsweise späterer Zeitpunkt. Lineare Systeme weisen eine endliche Ausgangsamplitude bei einer endlichen Eingangsamplitude auf. So 
    muss die Ausgangsamplitude mit dem Abfall des Systemeinganges ebenfalls enden.}
    \speech{SuperposLTI}{3}{Systeme die sowohl linear als auch zeitlich invariant sind, werden als eL Te Ii -Systeme bezeichnet, was die Grundlage dafür bildet,
    dass der Überlagerungssatz auf Gleichstromnetzwerke angewendet werden kann.}
\end{frame}

\begin{frame}
    \ftx{LTI-Systeme}

    \begin{Key point}{LTI-Systems}
        LTI systems represent linear and time-invariant systems. DC networks must be considered as LTI systems so that
        the superposition theorem can be applied to them. 
    \end{Key point}

    \speech{SuperposLTIMrk}{1}{LTI-Systeme stellen lineare und zeitinvariante Systeme dar. Die Gleich-stromnetzwerke müssen als LTI-Systeme betrachtet werden, 
    damit der Überlagerungssatz auf sie angewendet werden kann.}
\end{frame}



%Überlagerungssatz
\begin{frame}
    \renewcommand{\figurename}{Figure}
    \ftb{Superposition theorem}

    \s{
        Once the electrical networks have been checked for linearity as the systems under consideration, 
        the superposition theorem can be applied. Here, different sources can be considered separately from one another.
        In this way, the superposition theorem switches off all sources
        and switches on the individual sources in series. When sources are switched off, 
        current sources and voltage sources are converted differently. When an ideal
        voltage source is switched off, a short circuit occurs at the location of the voltage source, as shown in Figure \ref{BildSpannungErsetzen}. 
    }

    \b{
        Bedinungen und Ausführung des Überlagerungssatzes: 
        \begin{itemize}
            \item<1-> Vorausgesetzte Linearität
            \item<2-> Ausschalten der Quellen
            \item<3-> Ideale Spannungsquelle wird durch einen Kurzschluss ersetzt
            \item<4-> Ideale Stromquelle wird durch offene Klemmen, einen Leerlauf, ersetzt
        \end{itemize}
    }

    \onslide<3->{
    \fu{
        \resizebox{.6\textwidth}{!}{%
        \hspace{0.1cm}
        \begin{minipage}{0.3\textwidth}
            \input{Tikz/src/umwandlungspannungsquelle1.tex}	
		\end{minipage}
        \hspace{0.1cm}
        \begin{minipage}{0.2\textwidth}
            \input{Tikz/src/umwandlungspannungsquelle2.tex}	
		\end{minipage}
        \hspace{0.1cm}
        \begin{minipage}{0.3\textwidth}
            \input{Tikz/src/umwandlungspannungsquelle3.tex}		
		\end{minipage}
        \hspace{0.1cm}
        }
	}{{\bf Conversion of a voltage source.} If a voltage source is deactivated, a short circuit remains. This short circuit is used for further 
    network calculations.  \label{BildSpannungErsetzen}}}

    \s{
        When switched off, a power source leaves only open terminals (Figure \ref{BildStromErsetzen}). 
        This idle state prevents further consideration of this path in the network. 
    }

    \onslide<4->{
    \fu{
        \resizebox{.6\textwidth}{!}{%
        \hspace{0.1cm}
        \begin{minipage}{0.3\textwidth}
            \input{Tikz/src/umwandlungstromquelle1.tex}	
		\end{minipage}
        \hspace{0.1cm}
        \begin{minipage}{0.2\textwidth}
            \input{Tikz/src/umwandlungstromquelle2.tex}	
		\end{minipage}
        \hspace{0.1cm}
        \begin{minipage}{0.3\textwidth}
            \begin{circuitikz}
    \draw (0,0) to [short,-*] (0,1);
    \draw (0,3) to [short,-*] (0,2);
    \draw (1.5,1.5) node[right, align=center] {Idle state\\(open\\terminals)};
\end{circuitikz}		
		\end{minipage}
        \hspace{0.1cm}
        }
	}{{\bf Conversion of a power source.} Deactivating a power source leaves terminals open.   \label{BildStromErsetzen}}}

    
    \speech{SuperposQuell}{1}{Beim Überlagerungssatz wird ein elektrisches Netzwerk für jede Quelle berechnet. Hierfür wird die Linearität des Systems vorausgesetzt.}
    \speech{SuperposQuell}{2}{Um ein elektrisches Netzwerk für eine Quelle berechnen zu können, müssen die übrigen Quellen deaktiviert werden. Dabei gilt folgendes zu beachten:}
    \speech{SuperposQuell}{3}{Eine ideale Spannungsquelle erzeugt beim deaktivieren einen Kurzschluss.}
    \speech{SuperposQuell}{4}{Eine ideale Stromquelle hinterlässt nach der Deaktivierung offene Klemmen, welche keinen Stromfluss ermöglichen.}
\end{frame}

\begin{frame}
    \ftx{Conversion of sources}

    \begin{Merksatz}{Conversion of sources}
        When voltage sources and current sources are switched off, deactivated voltage sources leave a short circuit and 
        deactivated current sources leave an open circuit.
    \end{Merksatz}

    \speech{SuperposQuellMrk}{1}{Beim Ausschalten von Spannungsquellen und Stromquellen hinterlassen deaktivierte Spannungsquellen einen Kurzschluss und 
        deaktivierte Stromquellen einen Leerlauf.}
\end{frame}

\begin{frame}
    \ftx{Zusammenfassung Überlagerungssatz}

    \s{
        If all sources except one are switched off,
        the electrical network for the remaining source is analysed. This is then carried out sequentially 
        with each source. 
        At the end, the individual effects are considered as a sum. 
        For example, the current $I_\mathrm{R}$ through a resistor R, which is supplied by two sources $Q_1$ and $Q_2$,
        can be explained by the sum of the partial currents of the two sources (see Equation \ref{GleichungZweiQuellen}). 
    }

    \b{
        Annahme: 
        \begin{itemize}
            \item Der Widerstand $R$ wird durch die beiden Quellen $Q_1$ und $Q_2$ versorgt
            \item Separate Berechnung des Stromes durch $R$ 
            \item Gesamtstrom durch $R$ ergibt sich aus der Summe der Ströme aus den beiden Quellen:
        \end{itemize}
    }

    \begin{eq}
        I_\mathrm{R}=f(Q_1,Q_2)    \label{GleichungZweiQuellen}
    \end{eq}

    \speech{SuperposZsm}{1}{Angenommen ein Widerstand R wird durch die beiden Quellen Q 1 und Q 2 versorgt und der resultierende Strom durch den Widerstand 
    soll berechnet werden. Dann wird für jede Quelle der Strom durch den Widerstand bestimmt und hinterher der Gesamtstrom aus den beiden Teilströmen der 
    beiden Quellen errechnet.}
\end{frame}

\begin{frame}
    \renewcommand{\figurename}{Figure}
    \ftx{Beispiel}

    \s{
        To illustrate the superposition theorem, Figure \ref{BildÜberlagerungssatz} shows an electrical network with two 
        voltage sources and three resistors. To calculate this network with more than one source, 
        the network is considered separately for both voltage sources. 
    }

    \b{
        Berechnung eines Beispielnetzwerks:  

    }

    \fu{
        \resizebox{.6\textwidth}{!}{%
		\input{Tikz/src/ueberlagerungssatz.tex}
		}%
    }{{\bf Electrical network with two voltage sources to explain the superposition theorem.} The voltage sources should
    be analysed one after the other for the electrical network. \label{BildÜberlagerungssatz}}

    \s{
        The electrical network presented is shown again in Figure \ref{BildÜberlagerungssatzQuellen} for the 
        calculation of the two voltage sources separately. For the network analysis with voltage source
        $U_1$, voltage source $U_2$ is short-circuited. The two resistors $R_2$ and $R_3$ are now connected in parallel
        with each other. The voltage of $U_1$ is now distributed across $R_1$ and the parallel connection $R_{23}$. Equivalently,
        when considering the network for the voltage source $U_2$, the voltage source $U_1$ is short-circuited. Now
        the two resistors $R_1$ and $R_2$ form a parallel connection. The voltage of the voltage source $U_2$ is distributed across 
        the parallel connection $R_{12}$ and the resistor $R_3$.
    }

    \speech{SuperposAnw}{1}{Die zuvor erklärte Überlagerung wird nun anhand eines elektrischen Netzwerkes mit drei Widerständen und 
    zwei Spannungsquellen Berechnet. Es soll die Spannung über den Widerstand R 2 bestimmt werden.}
\end{frame}  



\begin{frame}
    \renewcommand{\figurename}{Figure}
    \ftx{Beispiel Lösung}

    \fu{
        \resizebox{0.8\textwidth}{!}{%
		\input{Tikz/src/ueberlagerungssatzquellen.tex}
		}%
    }{{\bf Electrical network with two voltage sources to explain the superposition theorem.} On the left, the network for analysing
    the voltage source $U_1$ is shown, and on the right, the network for analysing the voltage source $U_2$.  \label{BildÜberlagerungssatzQuellen}}

    \s{
        In equation \ref{GleichungU1} and equation \ref{GleichungU2}, the voltages for the two voltage sources
        $U_1$ and $U_2$ at resistor $R_3$ are determined separately. 
        In equation \ref{GleichungUR3}, the two voltages across the 
        resistor $R_3$ are then added together according to the superposition theorem.
    }
    
    
    \onslide<1->{
    \begin{eq}
        U_\mathrm{R3}(U_1) = U_1 \cdot \frac{R_\mathrm{23}}{R_1+R_\mathrm{23}} \label{GleichungU1} 
    \end{eq}}
    \onslide<2->{
    \begin{eq}
        U_\mathrm{R3}(U_2) = U_2 \cdot \frac{R_\mathrm{12}}{R_3+R_\mathrm{12}} \label{GleichungU2} 
    \end{eq}}
    \onslide<3->{
    \begin{eq}
        U_\mathrm{R3} = U_\mathrm{R3}(U_1) + U_\mathrm{R3}(U_2) \label{GleichungUR3}
    \end{eq}}

    \speech{SuperposAnwLsg}{1}{Zuerst wird das Netzwerk für die Berechnung mit der Spannungsquelle U eins vorbereitet und die Spannungsquelle U 2 deaktiviert.
    Die Widerstände R 2 und R 3 können zusammengefasst werden. Es ergibt sich die angezeigte Spannung über R 3 in der Abhängigkeit von U eins.}
    \speech{SuperposAnwLsg}{2}{Dasselbe wird für die Spannungsquelle U 2 durchgeführt.}
    \speech{SuperposAnwLsg}{3}{Am Ende werden beide Gleichungen zusammengerechnet.}
\end{frame}


\begin{frame}
    \ftx{Superposition principle}

    \begin{Merksatz}{Überlagerungssatz}
        Linear and time-invariant electrical networks with more than one source can be determined as the sum of the partial analyses of each 
        individual source. 
    \end{Merksatz}

    \speech{SuperposUberMrk}{1}{Lineare und zeitlich invariante elektrische Netzwerke mit mehr als einer Quelle können als Summe der Teilanalysen von jeder 
        einzelnen Quelle bestimmt werden. }
\end{frame}

\newpage
\begin{frame}
    \renewcommand{\figurename}{Figure}
    \ftx{Beispiel Superpositionsverfahren}
    \begin{bsp}{Superposition method}{}

    \s{
        The electrical network shown in Figure \ref{BildBeispielaufgabeSuperposition} is given. The following 
        tasks are to be completed: 
    }

        \begin{enumerate}
            \only<1>{	
                \item[]
                Application of the superposition principle (Superposition theorem): 
                \begin{enumerate}
                    \item[a)] Plotting currents and voltages
                    \item[b)] Network calculation $I_3$ for $U_\mathrm{g}$
                    \item[c)] Network calculation $I_3$ for $I_\mathrm{g}$
                    \item[d)] How big is the current? $I_3$
                \end{enumerate}

                
                \fu{
                    \input{Tikz/src/superposition1.tex}
                }{{\bf Example.} Network analysis using Kirchhoff's rules. \label{BildBeispielaufgabeSuperposition}}
                    }
			\only<2>{			
				\item[a)] Specification of currents and voltages:
                \begin{figure}[H]
                    \resizebox{0.65\textwidth}{!}{%   
                        \input{Tikz/src/superposition2.tex}
                    }\centering                                                           
                \end{figure}
                    }    
			\only<3>{
				\item[b)] Network calculation $I_3(U_\mathrm{g})$ for $U_\mathrm{g}$:
				\onslide<3>{
                    \begin{figure}[H]
                        \resizebox{0.65\textwidth}{!}{%  
                        \input{Tikz/src/superposition3.tex}
                        }\centering                                                           
                    \end{figure}
					\begin{eqa}
						K_1: I_\mathrm{ges}&=I_1+I_3 \nonumber\\                        
                        I_\mathrm{ges}&=\frac{U_\mathrm{g}}{R_\mathrm{ges}}=\frac{U_\mathrm{g}}{R_\mathrm{i}+\frac{R_1\cdot(R_2+R_3)}{R_1+R_2+R_3}} \nonumber\\
                        \frac{I_3}{I_\mathrm{ges}}&=\frac{R_1}{R_1+R_2+R_3} \nonumber\\
                        I_3(U_\mathrm{g})&=I_\mathrm{ges}\cdot\frac{R_1}{R_1+R_2+R_3}=\frac{U_\mathrm{g}}{R_\mathrm{i}+\frac{R_1\cdot(R_2+R_3)}{R_1+R_2+R_3}}\cdot\frac{R_1}{R_1+R_2+R_3}\nonumber
					\end{eqa}
				}
			}
			\only<4>{
				\item[c)] Network calculation $I_3(I_\mathrm{g})$ for $I_\mathrm{g}$:
				\onslide<4>{
                    \begin{figure}[H]
                        \resizebox{0.65\textwidth}{!}{%  
                        \input{Tikz/src/superposition4.tex}
                        }\centering                                                           
                    \end{figure}
					\begin{eqa}
						K_2: I_\mathrm{ges}&=I_\mathrm{g}=I_2-I_3 \nonumber\\                        
                        \frac{-I_3}{I_\mathrm{ges}}&=\frac{R_2}{(R_1||R_\mathrm{i})+R_2+R_3} \nonumber\\
                        I_3(I_\mathrm{g})&=-I_\mathrm{g}\cdot\frac{R_2}{(R_1||R_\mathrm{i})+R_2+R_3}\nonumber
				    \end{eqa}
				}
			}
            \only<5>{
				\item[d)] How strong is the current $I_3(U_\mathrm{g},I_\mathrm{g})$? \\ 

				\onslide<5>{
                    Superposition: \\
					\begin{eqa}
						I_3(U_\mathrm{g},I_\mathrm{g}) &= I_3(U_\mathrm{g})+I_3(I_\mathrm{g}) \nonumber\\
                        I_3(U_\mathrm{g},I_\mathrm{g}) &= \frac{U_\mathrm{g}}{R_\mathrm{i}+\frac{R_1\cdot(R_2+R_3)}{R_1+R_2+R_3}}\cdot\frac{R_1}{R_1+R_2+R_3}+(-I_\mathrm{g}\cdot\frac{R_2}{(R_1||R_\mathrm{i})+R_2+R_3}) \nonumber
				    \end{eqa}
				}
			}
		\end{enumerate}

    \end{bsp}

    \speech{SuperposBsp}{1}{Anhand des folgenden Beispieles sollen die Eklärungen zum Superpositionsprinzip noch einmal verfestigt werden.
    Hierzu sollen zu dem angezeigten elektrischen Netzwerk die Ströme und Spannungen eingezeichnet werden. 
    Der Strom I 3 durch den Widerstand R 3 soll getrennt in Abhängigkeit von der Spannungsquelle U Ge und der Stromquelle I Ge bestimmt werden.
    Der Gesamtstrom I 3 wird anschließend in Abhängigkeit beider Quellen zusammengefasst.}
    \speech{SuperposBsp}{2}{Hierzu werden die Spannungen und Ströme aller komponenten angegeben.}
    \speech{SuperposBsp}{3}{Für die Netzberechnung der Spannungsquelle U G wird die Stromquelle deaktiviert und hinterlässt die offenen Klemmen.
    Der Stromfluss über diesen Pfad ist somit unterbrochen. Der Gesamtstrom I ges teilt sich am Knoten K 1 auf die Ströme I 1 und I 3 auf. Der Strom 
    I 3 ist dabei gleich dem Strom I 2. Das Verhältnis von U G und dem Gesamtwiderstand R ges bestimmt I ges. Der Strom durch R 3 lässt sich hier über 
    das Verhältnis von I 3 zu I ges sowie R1 zu der Summe aus R 1, R 2 und R 3 ermitteln.}
    \speech{SuperposBsp}{4}{Bei der Netzbetrachtung für die Stromquelle I Ge wird die Spannungsquelle nun deaktiviert und hinterlässt einen Kurzschluss. 
    Der Gesamtstrom resultiert aus I Ge und teilt sich am Knoten Ka zwei auf die Ströme I zwei und minus I drei auf. 
    Nun lässt sich das Verhältnis zur Berechnung von I drei aufstellen. I drei verhält sich zu I gesamt wie R zwei zu der Summe aus der Parallelschaltung 
    von R eins und R i und R zwei und R drei.}
    \speech{SuperposBsp}{5}{Zum Schluss werden die Teilergebnisse aufaddiert.}
\end{frame}